\thispagestyle{otFront}
\phantomsection\pdfbookmark[0]{Références et repères éditoriaux}{bm-references}
\otPageTitre{Références}{France · Première générale · Spécialité mathématiques}
\fontsize{10.5}{14.5}\selectfont
\otIntertitre{Le programme de référence}
Le livre s'appuie sur le programme publié au \textbf{BO n° 14 du 2 avril 2026}
(arrêté du 26 février 2026, NOR MENE2602917A), applicable à la rentrée 2026-2027.
Les repères de l'annexe relient ses contenus aux cours et exercices du livre.
Les approfondissements facultatifs sont sélectionnés, sans prétendre les reproduire tous.\par\vspace{3mm}
\href{https://www.education.gouv.fr/bo/2026/Hebdo14/MENE2602917A}{\textcolor{otNavy}{Programme officiel et annexe sur education.gouv.fr}}\par\vspace{5mm}
\otIntertitre{L'épreuve anticipée}
La préparation suit la \textbf{note de service du 11 septembre 2026}, NOR MENE2622640N,
publiée au BO spécial n° 4 du 17 septembre 2026. Elle s'applique en 2026-2027 aux
épreuves présentées au titre du baccalauréat session 2028 et remplace la note de service
du 10 juin 2025 (NOR MENE2515469N). Les automatismes y prennent la forme de réponses
courtes. Le programme de Première précise que sa liste d'automatismes s'ajoute à celle
de Seconde, à entretenir en Première. Le guide du livre distingue ce cadre officiel du barème
choisi pour ses sujets d'entraînement.\par\vspace{3mm}
\href{https://www.education.gouv.fr/bo/2026/Special4/MENE2622640N}{\textcolor{otNavy}{Définition officielle de l'épreuve anticipée}}\par\vspace{2mm}
\href{https://eduscol.education.gouv.fr/5688/epreuve-anticipee-de-mathematiques-aux-baccalaureats-general-et-technologique}{\textcolor{otNavy}{Présentation et ressources de préparation sur Éduscol}}\par\vspace{5mm}
\otIntertitre{Les choix pédagogiques}
Les 210 exercices de chapitre, 14 défis et deux sujets blancs sont des créations pour ce livre.
Les sujets blancs ne reproduisent pas un sujet officiel. Les défis très exigeants sont
facultatifs et explicitent leurs prérequis ; ils ne nécessitent pas de cours de Terminale.
L'ordre du livre peut être adapté par le professeur.\par\vspace{3mm}
Les figures, courbes, arbres et tableaux sont vectoriels. Les schémas sans échelle
imposée sont signalés. Les programmes Python peuvent être exécutés séparément.
Les résultats numériques illustrent les calculs théoriques ; les arrondis et les
hypothèses d'indépendance sont précisés lorsqu'ils interviennent.\par\vspace{5mm}
{\fontsize{9}{12}\selectfont\color{otMuted}Sources officielles consultées le 23 octobre 2026.
La charte reprend le modèle OptiTAB de Terminale : couleurs, typographie et repères de difficulté.}
\clearpage
